\section{Fundamento teórico de la Entrega 4}
\label{sec:fundamento-teorico-e4}

Esta sección se separa deliberadamente de la Sección~\ref{sec:fundamento-teorico} (heredada de
las Entregas 1--3, Parte~\ref{part:heredado-e1-e3}): reúne el fundamento teórico específico de
los dos canales de acceso, el pipeline de integración continua y el protocolo de calidad de
producto que la Entrega 4 añade sobre la capa de datos distribuida ya existente.

\subsection{Arquitectura en capas y principios SOLID}

La arquitectura en capas separa un sistema en presentación, aplicación, dominio e
infraestructura~\cite{martin2017clean,evans2003ddd}, con una regla de dependencia única: el
código de una capa solo puede depender de capas más internas, nunca al revés —el dominio no
conoce ni la base de datos ni el protocolo de transporte que lo expone. Los principios SOLID son
las restricciones que impiden que esa separación se erosione con el tiempo: responsabilidad
única (una clase, un motivo de cambio), abierto/cerrado (extensible sin modificar código
existente), sustitución de Liskov (una subclase debe poder reemplazar a su clase base sin alterar
la corrección del programa), segregación de interfaces (interfaces pequeñas y específicas en vez
de una interfaz general que fuerza a implementar métodos irrelevantes), e inversión de
dependencias (los módulos de alto nivel no dependen de los de bajo nivel; ambos dependen de
abstracciones). El uso disciplinado de estos principios es prerrequisito para que los patrones de
diseño del catálogo de Gamma et al.~\cite{gamma1994patterns} aporten valor real, en vez de ser
indirección sin propósito.

\subsection{Ingeniería de aplicaciones móviles}

Las aplicaciones móviles enfrentan retos particulares frente al desarrollo de escritorio o web:
heterogeneidad de dispositivos, restricciones de batería y de red intermitente, un ciclo de vida
gestionado por el sistema operativo (no por el proceso de la propia aplicación), y ventanas de
publicación controladas por tiendas de aplicaciones~\cite{wasserman2010mobile}. La decisión entre
desarrollo nativo y multiplataforma no es una preferencia estética: tiene consecuencias medibles
en rendimiento, mantenibilidad y \emph{time-to-market}, documentadas empíricamente en la
literatura de ingeniería de software móvil~\cite{elkassas2017taxonomy,biornhansen2019crossplatform}.
La Sección~\ref{sec:movil} aplica dos de esos criterios cuantitativos —tamaño de artefacto y
dependencias de interoperabilidad— directamente sobre el proyecto construido, en vez de citarlos
solo como referencia teórica.

\subsection{Integración y entrega continuas}

La entrega continua se apoya en un \emph{pipeline} automatizado que ejecuta pruebas, empaqueta y
despliega en ambientes segregados~\cite{humble2010continuous}. La cultura DevOps subyacente
exige que ningún cambio llegue a un artefacto publicado sin pasar por los mismos
controles~\cite{kim2021devops}: en este proyecto, el pipeline de GitHub Actions
(Sección~\ref{sec:pruebas-cicd}) actúa como \emph{quality gate} explícito mediante dependencias
declaradas entre \emph{jobs} —si la integración falla, el despliegue de imágenes no ocurre.

\subsection{Observabilidad distribuida}

La observabilidad de un sistema distribuido se sustenta en tres señales: métricas (agregados
numéricos con etiquetas), registros estructurados, y trazas (secuencias de \emph{spans}
correlacionados por un identificador común)~\cite{sigelman2010dapper,beyer2016sre}. El estándar
OpenTelemetry~\cite{opentelemetry2024spec} unifica su modelo de datos y su exportación,
permitiendo que distintos proveedores de \emph{backend} de observabilidad (en este proyecto,
Grafana Tempo para trazas y Prometheus para métricas) consuman las mismas señales sin
instrumentación duplicada por proveedor.

\subsection{Modelo de calidad ISO/IEC 25010:2023}

El modelo ISO/IEC 25010:2023~\cite{iso25010} organiza la calidad de producto de software en ocho
características —adecuación funcional, eficiencia de desempeño, compatibilidad, usabilidad,
fiabilidad, seguridad, mantenibilidad y portabilidad—, cada una descompuesta en
subcaracterísticas y evaluable mediante métricas objetivas. Frente a un modelo de calidad
puramente cualitativo, esta estructura obliga a declarar de antemano un umbral numérico por
característica y a reportar la desviación respecto a ese umbral cuando ocurre, en vez de una
valoración subjetiva de si el sistema ``es de buena calidad''. La Sección~\ref{sec:iso25010}
aplica cinco de las ocho características a este proyecto, siguiendo esa misma disciplina de
declarar el umbral antes de medir.

\subsection{Trazabilidad de temas a artefacto}
\label{sec:trazabilidad-temas}

La Tabla~\ref{tab:trazabilidad-temas} hace explícito, para cada tema de la asignatura cubierto en
este documento, el artefacto real del repositorio donde se implementa —no solo el lugar donde se
menciona en el texto.

\begin{table}[H]
\centering
\caption{Trazabilidad de temas de la asignatura a su artefacto en el repositorio}
\label{tab:trazabilidad-temas}
\small
\begin{tabular}{@{}p{3.1cm}p{9.4cm}@{}}
\toprule
\textbf{Tema} & \textbf{Artefacto} \\
\midrule
gRPC &
  \texttt{services/telemetry-service/src/main/proto/telemetry.proto} y
  \texttt{TelemetryGrpcService.java} (servidor); cliente real en
  \texttt{ZonaVentanaTelemetriaStrategy} de \texttt{svc-principal} \\
Saga por coreografía sobre Kafka &
  \texttt{CreateTicketHandler} publica \texttt{TicketCreated}; \texttt{notification-service} y
  \texttt{report-service} (\texttt{ReportEventListener}) reaccionan de forma autónoma, sin
  orquestador central; \texttt{TechnicianEventPublisher} en \texttt{auth-service} sigue el mismo
  patrón para \texttt{TechnicianCreatedEvent} \\
CQRS &
  Separación por servicio: \texttt{svc-principal/application/command/} concentra la escritura
  (\texttt{CreateTicketHandler}, \texttt{AssignTechnicianHandler},
  \texttt{UpdateTicketStatusHandler}); \texttt{report-service/application/ReportQueryService.java}
  concentra la lectura, sobre una proyección propia (\texttt{TicketSummary}) \\
Fragmentación en CockroachDB &
  \texttt{db-cluster/scripts/init\_*.sql} (esquema particionado real); decisión documentada en
  \texttt{docs/adr/0003-sharding-policy.md} \\
Spark y ley de Amdahl &
  \texttt{spark/} (pipeline paralelo real); análisis en
  Sección~\ref{sec:pipeline-paralelo} y en la Subsección~\ref{sec:fundamento-modelo-paralelo} de
  la Sección~\ref{sec:fundamento-teorico} \\
Patrón Observer &
  \texttt{EscalationLoggingObserver}, \texttt{EscalationMetricsObserver} y
  \texttt{EscalationEventPublisherObserver} en \texttt{svc-principal}; decisión documentada en
  \texttt{docs/adr/0005-patrones-gof.md} \\
Arquitectura hexagonal (4 capas) &
  \texttt{domain/}, \texttt{application/}, \texttt{infrastructure/} y \texttt{presentation/} en
  \texttt{auth-service}, \texttt{report-service} y \texttt{svc-principal}; alcance real declarado
  en la Sección~\ref{sec:arquitectura-e4} \\
Observabilidad (Prometheus, Grafana, Tempo, OpenTelemetry) &
  \texttt{ops/prometheus/}, \texttt{ops/grafana/}, \texttt{ops/tempo/} y
  \texttt{ops/otel-collector/}; desarrollado en la Sección~\ref{sec:observabilidad} \\
\bottomrule
\end{tabular}
\end{table}
